% =====================================================================
%  01_introduccion.tex — Introducción, problemática, propuesta de
%  solución, objetivos y justificación.
% =====================================================================
\chapter{Introducción}
\label{cap:introduccion}

El Metrobús de la Ciudad de México opera con unidades que, al terminar o
iniciar su servicio, entran y salen de patios de encierro. Saber qué
unidades están dentro de un patio, desde cuándo y cuántas veces entraron o
salieron en un día es información básica para la operación: permite
planear la salida de unidades, detectar estancias anómalas y respaldar los
registros ante una disputa entre las partes.

En el Patio CCA este control se hacía de forma manual. El proyecto
\saru{} (\emph{Sistema Automatizado de Reconocimiento de Unidades}),
impulsado por Grupo CISA, busca automatizarlo con visión artificial: una
cámara en la puerta del patio, una computadora de borde que procesa el
video en tiempo real y un tablero web que muestra los movimientos al
momento en que ocurren~\cite{act1}.

\section{Problemática}
\label{sec:problematica}

El registro manual de entradas y salidas depende de que una persona anote
cada unidad que cruza la puerta. Esa dependencia genera tres problemas:

\begin{itemize}
  \item \textbf{Movimientos perdidos}, sobre todo en horas pico o cuando
        varias unidades cruzan al mismo tiempo.
  \item \textbf{Horas poco confiables}: la hora se anota tarde o se
        redondea, lo que altera el cálculo del tiempo que cada unidad
        permanece en el patio.
  \item \textbf{Falta de evidencia}: cuando un registro se cuestiona, no hay
        foto ni video que lo respalde.
\end{itemize}

Automatizar el registro tiene sus propias dificultades. El número
económico está pintado en distintas zonas de la unidad (frente, costado o
parte trasera) y la cámara no siempre lo ve completo; las unidades pueden
detenerse sobre la puerta, maniobrar o cruzar varias a la vez; y la red del
patio restringe la comunicación con servicios externos~\cite{rep-envivo}.
Un sistema útil debe resolver estos casos sin intervención humana
constante y con un costo de operación bajo.

\begin{aviso}[title={Por validar}]
Falta cuantificar el problema con datos del patio: cuántas unidades cruzan
por día, cuántos movimientos se perdían con el registro manual y quién
consume hoy esa información.
\end{aviso}

\section{Propuesta de solución}
\label{sec:propuesta}

Se propone un sistema de visión artificial que corre por completo en una
computadora NVIDIA Jetson Orin instalada junto a la puerta del patio, sin
depender de un servidor en la nube. A partir del video de la cámara, el
sistema:

\begin{enumerate}
  \item \textbf{Detecta} cada unidad con un modelo YOLO (M1) y la sigue
        cuadro a cuadro con ByteTrack.
  \item \textbf{Lee} el número económico en dos pasos: un segundo modelo
        (M2) localiza la zona donde está pintado y un tercero (M3) lee los
        dígitos. Un número solo se confirma cuando varias lecturas
        coinciden.
  \item \textbf{Decide} si la unidad entra o sale cuando su trayectoria
        atraviesa una línea virtual trazada sobre la imagen, y descarta las
        maniobras.
  \item \textbf{Registra} cada cruce y cada estancia en una base de datos
        SQLite, y compara el número con el padrón oficial de unidades.
  \item \textbf{Avisa} por Telegram y WhatsApp con hasta tres fotos del
        cruce; si la lectura es dudosa, la valida un modelo de lenguaje
        (Gemini) o una persona.
  \item \textbf{Publica} un tablero web con video en vivo, cruces, unidades
        dentro y estancias, y \textbf{graba} 72 horas de video continuo como
        evidencia.
\end{enumerate}

Los tres modelos se ejecutan en la GPU de la Jetson como motores TensorRT,
y el sistema se despliega en tres contenedores Docker: visión, grabación y
WhatsApp~\cite{rep-tecnico}.

\section{Objetivos}

\subsection{Objetivo general}

Desarrollar un sistema de visión artificial en el borde que registre de
forma automática, con evidencia y sin captura manual, cada entrada y salida
de unidades en la puerta del Patio CCA, y que ponga esa información a
disposición de los responsables en el momento en que ocurre.

\subsection{Objetivos específicos}

\begin{itemize}
  \item Detectar y seguir las unidades que aparecen en el video de la
        puerta, aunque haya varias a la vez.
  \item Leer el número económico de cada unidad en sus zonas delantera,
        lateral o trasera, y confirmarlo con varias lecturas.
  \item Determinar la dirección de cada cruce (entrada o salida) con una
        línea virtual, sin contar maniobras ni unidades detenidas.
  \item Registrar cruces y estancias en una base de datos, respetando el
        padrón oficial y conservando los números que no están en él.
  \item Notificar cada cruce con fotos y permitir su validación humana o
        automática.
  \item Ofrecer un tablero web para consultar los movimientos y conservar
        video continuo como evidencia.
  \item Validar el funcionamiento de cada prototipo con un protocolo de
        pruebas documentado.
\end{itemize}

\section{Justificación}

Un registro automático elimina la dependencia de una persona en la puerta,
sella la hora con el reloj del sistema y guarda evidencia de cada
movimiento: fotos del cruce y video continuo de las últimas 72 horas. Con
ello, las estancias, los tiempos y los conteos diarios del patio pasan a
ser datos confiables que pueden consultarse en cualquier momento.

La propuesta procesa todo \textbf{en el borde}. Esto reduce la dependencia
de la red del patio, que ya bloquea algunos servicios externos, evita
enviar video a terceros y mantiene un consumo bajo: entre 11 y 19~W según
haya o no unidades en cuadro~\cite{rep-tecnico}. El uso de tres modelos en
cascada (unidad, zona del número y dígitos) permite leer números pintados
en distintas partes de la unidad, algo que un detector único de matrículas
no resuelve.

En cuanto a viabilidad, el proyecto ya cuenta con el equipo instalado en el
patio, modelos entrenados para cada etapa y un sistema que opera en vivo.
El alcance de este documento comprende el primer prototipo funcional y la
planeación de los siguientes. El cierre del proyecto está previsto antes
del 30 de noviembre de 2026~\cite{rep-tecnico}.
